\section{研究背景}\label{sec:intro}
% !lp begin {"id":"intro-motivation","title":"研究动机","status":"final","tag":"终稿","note":"介绍背景与目标"}
写作研究报告时，作者常常在背景、方法、结果和结论之间切换。随着内容增长，找到需要修改的位置、记录各部分状态以及保持引用一致，都需要花费时间。

本示例以虚构的阅读习惯调查为主题，展示如何将文档组织为多个文件和片段。所有数字均为演示数据，不代表真实调查。
% !lp end
\subsection{研究问题}
% !lp begin {"id":"intro-questions","title":"研究问题","status":"draft","tag":"人工初稿","note":"补充研究假设"}
我们关注阅读时间如何分布，以及如何描述一段时间内的变化。方法将在第\ref{sec:method}节介绍。
\begin{enumerate}
\item 明确记录对象和观察周期；
\item 汇总每日记录，并检查缺失值；
\item 使用表格和公式呈现结果。
\end{enumerate}
% !lp end
\subsection{报告结构}
% !lp begin {"id":"intro-structure","title":"报告结构","status":"draft","tag":"AI初稿"}
本文从研究设计出发，随后展示示例数据，最后讨论结果的局限性。文献引用示例见\cite{knuth1984}。在导航台中可以分别打开章节和片段，修改后保存至原始文件。
% !lp end
